% Source: Pearce, Dooms, Rigg, Oramas, Sharkey,
% "Bilinear MLPs enable weight-based mechanistic interpretability,"
% ICLR 2025 (arXiv:2410.08417v2, 25 Jun 2025).
% Code: https://github.com/tdooms/bilinear-decomposition
%
% Factual corrections vs the user-supplied template (verified against the PDF):
%   1. The GLU sigmoid gate sits on V, not W: g(x) = (Wx) o sigma(Vx) (Eq. 1).
%      The template's "without the sigmoid gate" is correct; this comment is
%      just to flag the convention.
%   2. Equation numbers in the paper: the spectral decomposition Q = sum lambda_i v_i v_i^T
%      is Eq. 2; the quadratic form x^T Q x = sum lambda_i (v_i^T x)^2 is Eq. 3.
%      The template attributed the latter to "eq 3", which is correct.
%   3. Data-efficiency gap (Appendix I, Table 6): bilinear is 6% less data-efficient
%      than SwiGLU at constant epochs, but EQUALLY compute-efficient at constant
%      wall-clock time on a 6-layer / d_model=512 / d_hidden=2048 TinyStories
%      transformer. Reported as 1.337 vs 1.321 final loss (constant epochs);
%      1.337 vs 1.336 (constant time). Stated in paper as "SwiGLU attains the
%      same final loss of the bilinear variant in 6% fewer epochs."
%   4. TinyStories transformer in Section 5: 6-layer, ~33M params, SwiGLU
%      replaced by bilinear MLP (Appendix G.2). Template's "6-layer" is correct.
%   5. HOSVD (Section 3.3): the tensor is reshaped to (d_out, d_in^2) and a
%      standard SVD gives B = sum_i sigma_i u_out^{(i)} (x) q_{in x in}^{(i)},
%      where each q can be further eigendecomposed. Template's reshape is correct.
%   6. Limitations (Section 6): the method "typically relies on having a set of
%      meaningful output directions available. In shallow models, the unembedding
%      directions can be used, but in deeper models, we rely on features derived
%      from sparse autoencoders that are dependent on an input dataset."
%      Template's characterization (unembedding shallow / SAE deep, both
%      input-dataset-dependent) is correct.

% --- abstract sentence ---
% Insert immediately after the routed-CP framing claim in the abstract.
Whereas the closest prior weight-based interpretability programme for the
gated FFN~\citep{pearce2024bilinear} is restricted to bilinear MLPs --- a
substitute architecture in which the sigmoid gate of SwiGLU is removed so
that the layer becomes exactly a third-order tensor contracted twice with
the input --- the routed-CP form retains the gate as a per-channel router
and therefore applies directly to the SwiGLU layers used in production
language models.

% --- related-work paragraph ---
% Drop into the related-work / prior-art section.
\citet{pearce2024bilinear} study weight-based interpretability of bilinear
MLPs, $g(x) = (Wx) \odot (Vx)$, the GLU variant obtained by deleting the
sigmoid from $g(x) = (Wx) \odot \sigma(Vx)$ so that the layer is exactly
quadratic in the input; for a chosen output direction $u$ they contract the
third-order weight tensor $B$ to a symmetric interaction matrix
$Q = u \cdot_{\text{out}} B$ and read off the relevant input directions
from its eigendecomposition $Q = \sum_i \lambda_i v_i v_i^{\top}$, giving
$x^{\top} Q x = \sum_i \lambda_i (v_i^{\top} x)^2$ exactly (their Eq.~3).
This analysis is tight in the bilinear setting --- the $\alpha_j \equiv 1$
limit of our routed-CP form --- but does not lift to the sigmoid-gated
SwiGLU used in production models, where the identity
$\mathrm{SiLU}(z) = z \cdot \sigma(z)$ attaches an input-dependent
coefficient $\sigma(g_j^{\top} x)$ to every rank-one atom that the
eigendecomposition of a static $Q$ cannot see; consistent with this, they
report (Appendix~I, Table~6) that bilinear MLPs match SwiGLU at constant
compute but trail it by roughly $6\%$ in data efficiency on a $6$-layer
TinyStories transformer, and their language-model analysis requires output
directions supplied by an SAE trained on activations, reintroducing the
input-dataset dependence the weight-based programme aims to avoid. The
routed-CP form
$A(x) = \sum_j \alpha_j(x)\, u_j \otimes w_j \otimes g_j$ with
$\alpha_j(x) = \sigma(g_j^{\top} x)$ (Section~3) folds this gate into the
structural picture rather than removing it: each rank-one interaction atom
is selected by a learned per-channel router, and the static-bilinear
analysis of \citet{pearce2024bilinear} is recovered as the gate-saturated
special case. The two views are complementary --- their eigendecomposition
along an output direction surfaces the most important quadratic input
directions for that output, while CP/Tucker analysis along the channel
index $j$ surfaces the same-index coupling between $w_j$, $g_j$, and the
core slice $C^{(j)}$ that drives the separation result of Theorem~1.

% --- bibliography stub ---
% Add to your .bib (bibkey is provisional; rename as needed).
% @inproceedings{pearce2024bilinear,
%   title     = {Bilinear {MLP}s enable weight-based mechanistic interpretability},
%   author    = {Pearce, Michael T. and Dooms, Thomas and Rigg, Alice and
%                Oramas, Jose and Sharkey, Lee},
%   booktitle = {International Conference on Learning Representations (ICLR)},
%   year      = {2025},
%   eprint    = {2410.08417},
%   archivePrefix = {arXiv},
%   primaryClass  = {cs.LG},
%   url       = {https://arxiv.org/abs/2410.08417},
%   note      = {Code: \url{https://github.com/tdooms/bilinear-decomposition}}
% }
%
% Or, for a non-natbib bibliography:
% \bibitem{pearce2024bilinear}
%   M.~T. Pearce, T.~Dooms, A.~Rigg, J.~Oramas, and L.~Sharkey.
%   Bilinear {MLP}s enable weight-based mechanistic interpretability.
%   In \emph{International Conference on Learning Representations (ICLR)},
%   2025. arXiv:2410.08417.
